\PassOptionsToPackage{table,xcdraw}{xcolor}
\documentclass[11pt, a4paper]{article}

\usepackage[T1]{fontenc}
\usepackage[utf8]{inputenc}
\usepackage{tgpagella}
\usepackage[spanish, es-noshorthands]{babel}
\usepackage{amsmath, amssymb}
\usepackage{booktabs}
\usepackage{tabularx}
\usepackage{multicol}
\usepackage[most]{tcolorbox}
\usepackage[margin=2.5cm]{geometry}
\usepackage{fancyhdr}

\pagestyle{fancy}
\fancyhf{}
\setlength{\headheight}{16pt}
\lhead{\textbf{UCB Sede Tarija} | IMT-221}
\chead{Estructura Base: Informe IEEE}
\rhead{Circuitos Electrónicos II}
\renewcommand{\headrulewidth}{0.4pt}
\definecolor{ucbblue}{RGB}{0, 51, 102}

\begin{document}

\begin{center}
    \Large\textbf{\textcolor{ucbblue}{Estructura de Informe de Laboratorio — Tipo IEEE}}[cite: 17]
\end{center}

\textit{Nota: Esta plantilla adapta la estructura de un artículo IEEE a un informe académico. La maquetación final debe realizarse usando la plantilla IEEE oficial en Word o LaTeX[cite: 17].}

\section*{Título, Autores, Abstract y Keywords[cite: 17]}
\textbf{Título:} Específico y descriptivo. \textit{Ej: Caracterización DC y Compensación de Offset de un Amplificador Diferencial con LM741[cite: 17].} \\
\textbf{Abstract (100–180 palabras):} Resuma en un solo párrafo qué problema se estudió, el circuito usado, el procedimiento, el resultado cuantitativo principal y la conclusión técnica[cite: 17]. No lo convierta en una introducción[cite: 17]. \\
\textbf{Keywords:} 3 a 5 términos clave (ej. \textit{operational amplifier, differential amplifier, slew rate})[cite: 17].

\section*{I. Introduction[cite: 17]}
Responda: ¿Qué fenómeno se estudia y por qué importa en circuitos reales?[cite: 17] ¿Qué pregunta experimental aborda el laboratorio?[cite: 17] Finalice con un objetivo técnico claro, sin incluir los resultados completos[cite: 17].

\section*{II. Theoretical Background[cite: 17]}
Incluya \textbf{únicamente} la teoría necesaria para interpretar el experimento[cite: 17].
\begin{itemize}
    \item \textbf{A. Circuit Model:} Defina $V_1, V_2, V_o$ y escriba la ecuación de transferencia del circuito bajo prueba[cite: 17].
    \item \textbf{B. Real Device Effects:} Conceptos relevantes (ej. $V_{OS}$, $I_B$, mismatch resistivo, GBW, SR)[cite: 17]. Toda especificación real debe provenir de un datasheet[cite: 17].
\end{itemize}

\section*{III. Methodology[cite: 17]}
\begin{itemize}
    \item \textbf{A. Components and Instruments:} Use una tabla detallando modelo, valor y cantidad[cite: 17].
    \item \textbf{B. Circuit:} Esquema eléctrico con nombres de nodos, valores, alimentación, y referencias[cite: 17].
    \item \textbf{C. Experimental Procedure:} Describa la secuencia \textbf{real} ejecutada (ej. verificación de pinout, medición de rieles, troubleshooting, compensación). No copie la guía del docente[cite: 17].
\end{itemize}

\section*{IV. Results[cite: 17]}
\begin{itemize}
    \item \textbf{A. Measurements:} Tabule condiciones de prueba, voltajes ideales, medidos y errores absolutos[cite: 17].
    \item \textbf{B. Before / After Compensation:} Tabule la variación de $V_o$ y el error absoluto[cite: 17].
    \item \textbf{C. Figures:} Inserte capturas de osciloscopio, Bode plots o gráficas comparativas. Toda figura debe estar numerada, tener un título, estar referenciada en el texto y poseer unidades[cite: 17].
\end{itemize}

\section*{V. Discussion[cite: 17]}
\textbf{Sección crítica del informe.[cite: 17]} No repita resultados, interprételos[cite: 17]. Responda: ¿Coincidió el experimento con la predicción?[cite: 17] ¿Qué discrepancia apareció y qué mecanismos la explican?[cite: 17] ¿Qué evidencia apoya cada hipótesis?[cite: 17] ¿Qué cambió tras la corrección y qué error permaneció?[cite: 17] Distinga entre observación, interpretación y conclusión[cite: 17].

\section*{VI. Conclusion \& References[cite: 17]}
\textbf{Conclusion:} Responda al objetivo. Incluya el resultado técnico, evidencia cuantitativa y significado[cite: 17]. Evite "el laboratorio fue exitoso"[cite: 17]. \\
\textbf{References:} Numeradas por orden de aparición[cite: 17]. Cite datasheets correctamente y evite webs genéricas[cite: 17].

\begin{tcolorbox}[colback=white, colframe=black, title=\textbf{Comparación Recomendada para el Informe[cite: 17]}]
    \centering $\boxed{ \text{teoría} \leftrightarrow \text{simulación} \leftrightarrow \text{experimento} }$[cite: 17]
\end{tcolorbox}

\end{document}
